% GENERATED FILE. Edit canonical lessons, then run npm run generate:books.
% canonical-source-hash: fnv1a64:efe145f7096cc8a6
% canonical-lessons: FR-C390-budget, FR-C390-revenu, FR-C390-abonnement, FR-C390-depense, FR-C390-cotisation

\chapter{Les comptes de la maison}
\label{ch:fr-les-comptes-de-la-maison}
\hlchaptermodality{390}

\begin{chapteropening}
\textbf{By the end of this chapter:} \emph{I can talk about household money: the budget, income, a subscription, spending and a membership fee.}

Complete the last lesson of chapter 390: I can talk about household money: the budget, income, a subscription, spending and a membership fee.
\end{chapteropening}

\section[le budget]{le budget — the budget}

\label{lesson:FR-C390-budget}



\begin{quote}
Before the new one: say the French for to testify, then the French for the victim.
\end{quote}

\subsection*{You'll want to know: le budget}
\textbf{le budget} — \textquotedblleft{}the budget\textquotedblright{}.

\begin{quote}
Notre budget pour les vacances est de mille euros.
\end{quote}

\begin{quote}
\emph{Our holiday budget is a thousand euros.}
\end{quote}

The final \emph{t} is silent.

\begin{cousinweb}
\emph{budget} was borrowed from English \emph{budget} (EtymDB), which came from Old French \emph{bougette} \textquotedblleft{}a little leather bag\textquotedblright{}, a small form of \emph{bouge} \textquotedblleft{}bag\textquotedblright{}. \emph{bouge} goes back to Late Latin \emph{bulga} \textquotedblleft{}a leather bag\textquotedblright{}, a Gaulish word (EtymDB; Lewis \& Short).
\end{cousinweb}

\begin{grammarlens}[title={What you've built}]
One more word for telling, explaining and describing at B1.
\end{grammarlens}

\subsection*{Your turn}
\begin{itemize}
\raggedright
  \item \emph{Say it:} \emph{le budget}
  \item \emph{Say it:} \emph{le budget}, once more
  \item \emph{Say it:} a sentence of your own with \emph{le budget}
  \item \emph{Recall:} say the French for to testify, then the French for the victim, then say \emph{le budget} again
\end{itemize}

\begin{tcolorbox}[breakable,colback=teal!4,colframe=teal!35!black,title={Before you move on}]
What does le budget mean? (\textquotedblleft{}The budget\textquotedblright{}.) Say it once more.
\end{tcolorbox}
\section[le revenu]{le revenu — the income}

\label{lesson:FR-C390-revenu}



\begin{quote}
Before the new one: say the French for the victim, then the French for the budget.
\end{quote}

\subsection*{You'll want to know: le revenu}
\textbf{le revenu} — \textquotedblleft{}the income\textquotedblright{}.

\begin{quote}
Leur revenu a baissé cette année.
\end{quote}

\begin{quote}
\emph{Their income went down this year.}
\end{quote}

Often plural: \emph{les revenus}. From \emph{revenir}: what comes back to you.

\begin{cousinweb}
\emph{revenu} is the past participle of \emph{revenir} \textquotedblleft{}to come back\textquotedblright{} used as a noun: what comes back to you. \emph{revenir} comes from Latin \emph{revenīre}, built on \emph{venīre} \textquotedblleft{}to come\textquotedblright{} (Lewis \& Short).
\end{cousinweb}

\begin{grammarlens}[title={What you've built}]
One more word for telling, explaining and describing at B1.
\end{grammarlens}

\subsection*{Your turn}
\begin{itemize}
\raggedright
  \item \emph{Say it:} \emph{le revenu}
  \item \emph{Say it:} \emph{le revenu}, once more
  \item \emph{Say it:} a sentence of your own with \emph{le revenu}
  \item \emph{Recall:} say the French for the victim, then the French for the budget, then say \emph{le revenu} again
\end{itemize}

\begin{tcolorbox}[breakable,colback=teal!4,colframe=teal!35!black,title={Before you move on}]
What does le revenu mean? (\textquotedblleft{}The income\textquotedblright{}.) Say it once more.
\end{tcolorbox}
\section[l'abonnement]{l'abonnement — the subscription; the season ticket}

\label{lesson:FR-C390-abonnement}



\begin{quote}
Before the new one: say the French for the budget, then the French for the income.
\end{quote}

\subsection*{You'll want to know: l'abonnement}
\textbf{l'abonnement} — \textquotedblleft{}the subscription; the season ticket\textquotedblright{}.

\begin{quote}
J'ai un abonnement de bus pour l'année.
\end{quote}

\begin{quote}
\emph{I have a bus pass for the year.}
\end{quote}

\emph{S'abonner à} — to subscribe to. Masculine.

\begin{cousinweb}
\emph{abonnement} is built on \emph{abonner} \textquotedblleft{}to subscribe\textquotedblright{} (EtymDB), usually traced to Old French \emph{bonne} \textquotedblleft{}a boundary, a limit\textquotedblright{}; first, fixing a set sum in advance.
\end{cousinweb}

\begin{grammarlens}[title={What you've built}]
One more word for telling, explaining and describing at B1.
\end{grammarlens}

\subsection*{Your turn}
\begin{itemize}
\raggedright
  \item \emph{Say it:} \emph{l'abonnement}
  \item \emph{Say it:} \emph{l'abonnement}, once more
  \item \emph{Say it:} a sentence of your own with \emph{l'abonnement}
  \item \emph{Recall:} say the French for the budget, then the French for the income, then say \emph{l'abonnement} again
\end{itemize}

\begin{tcolorbox}[breakable,colback=teal!4,colframe=teal!35!black,title={Before you move on}]
What does l'abonnement mean? (\textquotedblleft{}The subscription; the season ticket\textquotedblright{}.) Say it once more.
\end{tcolorbox}
\section[la dépense]{la dépense — the expense, the spending}

\label{lesson:FR-C390-depense}



\begin{quote}
Before the new one: say the French for the income, then the French for the subscription.
\end{quote}

\subsection*{You'll want to know: la dépense}
\textbf{la dépense} — \textquotedblleft{}the expense, the spending\textquotedblright{}.

\begin{quote}
Le loyer est notre plus grosse dépense.
\end{quote}

\begin{quote}
\emph{Rent is our biggest expense.}
\end{quote}

From \emph{dépenser}.

\begin{cousinweb}
\emph{dépense} comes through Middle French \emph{despense} (EtymDB) from Late Latin \emph{dispēnsa}, from \emph{dispendere} \textquotedblleft{}to weigh out, to pay out\textquotedblright{}, built on \emph{pendere} \textquotedblleft{}to weigh\textquotedblright{} (Lewis \& Short).
\end{cousinweb}

\begin{grammarlens}[title={What you've built}]
One more word for telling, explaining and describing at B1.
\end{grammarlens}

\subsection*{Your turn}
\begin{itemize}
\raggedright
  \item \emph{Say it:} \emph{la dépense}
  \item \emph{Say it:} \emph{la dépense}, once more
  \item \emph{Say it:} a sentence of your own with \emph{la dépense}
  \item \emph{Recall:} say the French for the income, then the French for the subscription, then say \emph{la dépense} again
\end{itemize}

\begin{tcolorbox}[breakable,colback=teal!4,colframe=teal!35!black,title={Before you move on}]
What does la dépense mean? (\textquotedblleft{}The expense, the spending\textquotedblright{}.) Say it once more.
\end{tcolorbox}
\section[la cotisation]{la cotisation — the contribution, the membership fee}

\label{lesson:FR-C390-cotisation}



\begin{quote}
Before the new one: say the French for the subscription, then the French for the expense.
\end{quote}

\subsection*{You'll want to know: la cotisation}
\textbf{la cotisation} — \textquotedblleft{}the contribution, the membership fee\textquotedblright{}.

\begin{quote}
La cotisation au club de tennis est de cent euros.
\end{quote}

\begin{quote}
\emph{The tennis club fee is a hundred euros.}
\end{quote}

On a payslip, \emph{les cotisations} are social security contributions.

\begin{cousinweb}
\emph{cotisation} is built on \emph{cotiser} \textquotedblleft{}to pay one's share\textquotedblright{} (EtymDB), from \emph{cote} \textquotedblleft{}a share\textquotedblright{}, from Medieval Latin \emph{quota} \textquotedblleft{}how great a part\textquotedblright{}, from \emph{quotus} \textquotedblleft{}how many\textquotedblright{} (Lewis \& Short).
\end{cousinweb}

\begin{grammarlens}[title={What you've built}]
One more word for telling, explaining and describing at B1.
\end{grammarlens}

\subsection*{Your turn}
\begin{itemize}
\raggedright
  \item \emph{Say it:} \emph{la cotisation}
  \item \emph{Say it:} \emph{la cotisation}, once more
  \item \emph{Say it:} a sentence of your own with \emph{la cotisation}
  \item \emph{Recall:} say the French for the subscription, then the French for the expense, then say \emph{la cotisation} again
\end{itemize}

\begin{tcolorbox}[breakable,colback=teal!4,colframe=teal!35!black,title={Before you move on}]
What does la cotisation mean? (\textquotedblleft{}The contribution, the membership fee\textquotedblright{}.) Say it once more.
\end{tcolorbox}
